\section{Discussion}

The comprehensive hyperparameter optimization conducted via the Optuna framework, across multiple disparate datasets, provides a robust confirmation of the working hypothesis. The results suggest a paradigm shift in EEG interpolation: graph-aware and temporal-regularization (TV) methods are not merely theoretically competitive, but practically superior in realistic missing-channel scenarios when properly tuned.

\subsection{Physiological and Temporal Fidelity (ERP \& PSD)}
A critical finding from the optimization phase is the profound impact of temporal regularization on physiological signal features. Even when geometric baselines (e.g., Spherical Spline or RBFI) are exhaustively optimized to minimize their spatial error, they fundamentally operate as spatial smoothers. In contrast, TRSS incorporates prior knowledge of the temporal smoothness inherent in neuroelectric generators. 

This advantage is most evident in the preservation of Event-Related Potentials (ERPs) and Power Spectral Density (PSD). As shown previously in the Results section (\figurename~\ref{fig:erp_optimized} and \figurename~\ref{fig:psd_optimized}), TRSS accurately recovers the amplitude and latency of ERP components (such as the P300 or N100) even under severe contiguous channel loss (40\% nearby missing channels). Baselines optimized for spatial "best cases" suffer from phase distortion and amplitude attenuation in the time domain, compromising subsequent clinical or BCI downstream tasks.

\subsection{Multi-Dataset Generalization}
The variance in performance across datasets underscores the importance of the Optuna optimization framework. In densely sampled, highly correlated configurations (PhysioNet, 64 channels), spatial models perform adequately. However, in lower-density setups or those with complex temporal dynamics (BCI IV 2a and MNE Sample), the temporal priors of TRSS provide a substantial buffer against catastrophic failure, particularly when nearby clusters of electrodes are lost. 

\subsection{Concluding Remarks on Optimization}
Historically, novel algorithms are often compared against default-configured baselines, inflating perceived improvements. By subjecting all methods---including classical splines---to rigorous, dataset-specific hyperparameter search, this study establishes a fair baseline. The fact that TRSS and Tikhonov regularization maintain statistically significant margins over optimized baselines validates the fundamental premise: integrating spatial graph topology with temporal smoothness represents the state-of-the-art approach for robust EEG signal recovery.
